\documentclass[12pt,a4paper]{article}

\usepackage[utf8]{inputenc}
\usepackage[T1]{fontenc}
\usepackage[brazil]{babel}
\usepackage{geometry}
\usepackage{graphicx}
\usepackage{hyperref}
\usepackage{longtable}
\usepackage{array}
\usepackage{xcolor}
\usepackage{enumitem}
\usepackage{float}

\geometry{margin=2.5cm}
\hypersetup{
    colorlinks=true,
    linkcolor=blue,
    urlcolor=blue,
    citecolor=blue
}

\newcommand{\screenshot}[2]{
    \begin{figure}[H]
        \centering
        \fbox{\parbox[c][6cm][c]{0.9\textwidth}{\centering Inserir screenshot: #1}}
        \caption{#2}
    \end{figure}
}

\title{Manual de Utilização do Quintana}
\author{Equipe do Projeto}
\date{\today}

\begin{document}

\maketitle

\begin{center}
\textbf{Versão do documento:} preencher \\
\textbf{Versão/commit do sistema:} preencher \\
\textbf{Responsáveis pela documentação:} preencher
\end{center}

\newpage
\tableofcontents
\newpage

\section{Apresentação}

\subsection{Objetivo deste manual}
Explicar que este documento apresenta o passo a passo de utilização do Quintana por estudantes e professores.

\subsection{Público-alvo}
Indicar os públicos:
\begin{itemize}
    \item estudantes;
    \item professores;
    \item mediadores de oficinas;
    \item equipe de apoio.
\end{itemize}

\subsection{O que este manual não cobre}
Indicar que instalação, arquitetura, banco de dados e detalhes técnicos ficam no manual do sistema ou na documentação técnica.

\subsection{Convenções}
Explicar como serão representados:
\begin{itemize}
    \item nomes de botões;
    \item nomes de abas;
    \item mensagens do sistema;
    \item exemplos de usuários;
    \item screenshots.
\end{itemize}

\section{Acesso ao sistema}

\subsection{Endereço da aplicação}
Informar a URL de acesso usada na oficina ou ambiente alvo.

\screenshot{Tela de login}{Tela inicial de login do Quintana}

\subsection{Entrar no sistema}
Passo a passo:
\begin{enumerate}
    \item acessar a URL da aplicação;
    \item informar e-mail;
    \item informar senha;
    \item clicar em \textbf{Entrar}.
\end{enumerate}

\subsection{Sair do sistema}
Passo a passo para encerrar a sessão.

\subsection{Usuários de demonstração}

\begin{longtable}{|p{3cm}|p{5cm}|p{4cm}|}
\hline
\textbf{Perfil} & \textbf{E-mail} & \textbf{Senha} \\
\hline
Professor & preencher & preencher \\
\hline
Aluno & preencher & preencher \\
\hline
\end{longtable}

\section{Utilização pelo estudante}

\subsection{Visão geral da área do estudante}
Descrever a primeira tela após login como estudante.

\screenshot{Home do estudante}{Página inicial do estudante}

\subsection{Consultar temas disponíveis}
Passo a passo para localizar um tema.

\screenshot{Aba Temas}{Lista de temas disponíveis para o estudante}

\subsection{Enviar nova redação}
Passo a passo:
\begin{enumerate}
    \item escolher um tema;
    \item clicar em \textbf{Inserir Nova Redação};
    \item escrever ou colar o texto;
    \item enviar a redação;
    \item aguardar a correção.
\end{enumerate}

\screenshot{Tela de envio de redação}{Tela usada pelo estudante para escrever e enviar uma redação}

\subsection{Consultar minhas redações}
Explicar a aba \textbf{Minhas redações}.

\screenshot{Aba Minhas redações}{Lista de redações enviadas pelo estudante}

\subsection{Interpretar os cards de resumo}
Explicar os cards apresentados no topo da página:
\begin{itemize}
    \item Redações;
    \item Temas disponíveis;
    \item Última nota;
    \item Prioridade atual.
\end{itemize}

\subsection{Interpretar o radar médio das competências}
Explicar:
\begin{itemize}
    \item o que o radar mostra;
    \item o significado das competências C1 a C5;
    \item como identificar pontos fortes;
    \item como identificar prioridades de estudo.
\end{itemize}

\screenshot{Radar médio das competências}{Radar médio exibido na aba Minhas redações}

\subsection{Interpretar a linha do tempo de progresso}
Explicar:
\begin{itemize}
    \item evolução da nota total;
    \item evolução por competência;
    \item maior evolução;
    \item competência mais estável;
    \item queda recente.
\end{itemize}

\screenshot{Linha do tempo}{Linha do tempo de progresso do estudante}

\subsection{Abrir detalhes de uma redação}
Passo a passo para abrir o modal de detalhes.

\screenshot{Modal de redação}{Detalhes de uma redação}

\subsection{Aba Resumo}
Explicar os elementos da aba \textbf{Resumo}:
\begin{itemize}
    \item nota total;
    \item radar individual da redação;
    \item prioridades de estudo.
\end{itemize}

\subsection{Aba Texto}
Explicar onde o estudante consulta o texto enviado.

\subsection{Aba Notas}
Explicar as notas por competência.

\subsection{Aba Plano de ação}
Explicar:
\begin{itemize}
    \item mapa de feedback por competência;
    \item diagnóstico;
    \item sugestão de melhoria;
    \item ação prática;
    \item checklist de reescrita.
\end{itemize}

\screenshot{Plano de ação}{Mapa de feedback e checklist de reescrita}

\subsection{Aba Versões}
Explicar como visualizar comparação entre versões da mesma redação.

\subsection{Aba Feedback}
Explicar onde aparece o feedback textual.

\subsection{Reescrever uma redação}
Passo a passo:
\begin{enumerate}
    \item abrir uma redação;
    \item clicar em \textbf{Reescrever};
    \item produzir nova versão;
    \item enviar;
    \item comparar versões.
\end{enumerate}

\subsection{Consultar atividades atribuídas}
Explicar a aba \textbf{Atividades}.

\screenshot{Aba Atividades}{Atividades atribuídas ao estudante}

\subsection{Entender status das atividades}
Explicar:
\begin{itemize}
    \item pendente;
    \item enviada;
    \item atrasada.
\end{itemize}

\section{Utilização pelo professor}

\subsection{Visão geral da área do professor}
Descrever a primeira tela após login como professor.

\screenshot{Home do professor}{Página inicial do professor}

\subsection{Interpretar os cards de resumo}
Explicar os cards:
\begin{itemize}
    \item Redações;
    \item Temas disponíveis;
    \item Última nota;
    \item Alunos listados.
\end{itemize}

\subsection{Gerenciar temas}
Explicar a aba \textbf{Temas}.

\screenshot{Aba Temas do professor}{Lista de temas do professor}

\subsubsection{Criar novo tema}
Passo a passo:
\begin{enumerate}
    \item clicar em \textbf{Adicionar Tema};
    \item preencher nome do tema;
    \item preencher descrição;
    \item salvar.
\end{enumerate}

\subsubsection{Editar tema}
Passo a passo para editar um tema existente.

\subsubsection{Remover tema}
Passo a passo para remover um tema.

\subsection{Consultar redações recebidas}
Explicar a aba \textbf{Redações recebidas}.

\screenshot{Redações recebidas}{Tabela de redações recebidas pelo professor}

\subsection{Filtrar redações}
Explicar filtros disponíveis:
\begin{itemize}
    \item todas as redações;
    \item redações dos meus temas;
    \item aluno.
\end{itemize}

\subsection{Abrir e corrigir uma redação}
Passo a passo:
\begin{enumerate}
    \item clicar no título da redação;
    \item consultar texto e notas do modelo;
    \item preencher notas do professor, quando aplicável;
    \item preencher feedback do professor;
    \item salvar.
\end{enumerate}

\screenshot{Correção do professor}{Tela de edição ou correção de redação pelo professor}

\subsection{Analisar a turma}
Explicar a aba \textbf{Análise da turma}.

\screenshot{Análise da turma}{Painel de analytics do professor}

\subsection{Filtros da análise}
Explicar:
\begin{itemize}
    \item filtro por turma;
    \item filtro por atividade;
    \item agrupamento por atividade, tema ou dia.
\end{itemize}

\subsection{Distribuição por competência}
Explicar como interpretar a distribuição das notas da turma.

\subsection{Heatmap aluno por competência}
Explicar como identificar padrões no heatmap.

\subsection{Ranking de competências problemáticas}
Explicar como usar o ranking para planejamento docente.

\subsection{Agrupamento pedagógico de alunos}
Explicar como interpretar os grupos sugeridos.

\subsection{Evolução da turma}
Explicar como interpretar a evolução ao longo de atividades, temas ou datas.

\subsection{Comparação entre tema e desempenho}
Explicar como identificar temas com maior dificuldade.

\subsection{Alertas pedagógicos}
Explicar os tipos de alerta e como eles podem apoiar intervenção docente.

\subsection{Gerenciar turmas e atividades}
Explicar a aba \textbf{Turmas e atividades}.

\screenshot{Turmas e atividades}{Tela de gerenciamento de turmas e atividades}

\subsubsection{Criar turma}
Passo a passo para criar uma turma e selecionar estudantes.

\subsubsection{Editar turma}
Passo a passo para editar nome ou estudantes da turma.

\subsubsection{Remover turma}
Passo a passo para remover turma.

\subsubsection{Criar atividade}
Passo a passo para criar atividade vinculada a turma e tema.

\subsubsection{Editar atividade}
Passo a passo para editar atividade.

\subsubsection{Remover atividade}
Passo a passo para remover atividade.

\subsection{Acompanhar submissões de atividade}
Explicar como consultar:
\begin{itemize}
    \item alunos esperados;
    \item alunos que enviaram;
    \item alunos pendentes;
    \item alunos atrasados.
\end{itemize}

\section{Utilização em oficinas}

\subsection{Roteiro sugerido para estudantes}
Inserir roteiro curto de uso durante oficina.

\subsection{Roteiro sugerido para professores}
Inserir roteiro curto de uso durante oficina.

\subsection{Atividades práticas sugeridas}
Listar sugestões de atividades para oficina.

\subsection{Dicas para mediação}
Indicar cuidados ao apresentar notas, feedback e analytics.

\section{Mensagens comuns e solução de problemas}

\subsection{Não consigo entrar}
Orientar verificações de e-mail, senha e backend.

\subsection{A página mostra 404}
Informar rotas corretas, como `/quintana/login` e `/quintana/home`.

\subsection{Minhas redações não aparecem}
Orientar recarregar, conferir login e chamar equipe de apoio.

\subsection{Analytics não carrega}
Orientar verificar filtros, turmas, atividades e conexão.

\subsection{Sessão expirada}
Explicar que o usuário deve fazer login novamente.

\subsection{Acesso negado}
Explicar que o usuário pode estar tentando acessar recurso de outro perfil.

\section{Perguntas frequentes}

\subsection{A nota automática substitui o professor?}
Responder.

\subsection{O que significam C1, C2, C3, C4 e C5?}
Responder.

\subsection{Posso reescrever a mesma redação?}
Responder.

\subsection{O professor vê minhas redações?}
Responder.

\subsection{O que faço se minha nota baixar?}
Responder.

\section{Glossário rápido}

\begin{longtable}{|p{4cm}|p{10cm}|}
\hline
\textbf{Termo} & \textbf{Definição} \\
\hline
Competência & preencher \\
\hline
Tema & preencher \\
\hline
Redação & preencher \\
\hline
Reescrita & preencher \\
\hline
Atividade & preencher \\
\hline
Turma & preencher \\
\hline
Radar & preencher \\
\hline
Heatmap & preencher \\
\hline
\end{longtable}

\section{Apêndices}

\subsection{Apêndice A: checklist rápido do estudante}
Inserir checklist de ações principais do estudante.

\subsection{Apêndice B: checklist rápido do professor}
Inserir checklist de ações principais do professor.

\subsection{Apêndice C: usuários de demonstração}
Preencher usuários e senhas usados na oficina.

\subsection{Apêndice D: histórico de versões do manual}

\begin{longtable}{|p{3cm}|p{4cm}|p{7cm}|}
\hline
\textbf{Data} & \textbf{Responsável} & \textbf{Alteração} \\
\hline
preencher & preencher & preencher \\
\hline
\end{longtable}

\end{document}
